\subsection{用户展示模块}

\paragraph*{业务流程}

\begin{center}
\resizebox{0.98\linewidth}{!}{\begin{tikzpicture}[
  node distance=8mm,
  box/.style={draw,rounded corners,align=center,minimum height=9mm,inner sep=4pt,font=\small},
  arr/.style={-{Stealth[length=2mm]},semithick}
]
  \node[box] (login) {登录};
  \node[box,right=of login] (center) {个人中心};
  \node[box,right=of center] (records) {检测记录};
  \node[box,below=of records] (detail) {任务详情};
  \node[box,left=of detail] (report) {综合鉴伪报告};
  \node[box,left=of report] (audit) {人工审核进度};
  \draw[arr] (login) -- (center) -- (records);
  \draw[arr] (records) -- (detail);
  \draw[arr] (detail) -- (report) -- (audit);
\end{tikzpicture}
}
\end{center}

\paragraph*{用例图}

\begin{center}
\resizebox{0.98\linewidth}{!}{% 用户展示模块整体用例图（FR-YHZS-0001～0006，与 03-2-user-display.tex 一致）
\begin{tikzpicture}[
  uc/.style={draw,ellipse,minimum width=2.85cm,minimum height=9mm,align=center,font=\scriptsize,inner sep=2pt},
  arr/.style={semithick}
]
  \node[draw,rounded corners,minimum width=8.8cm,minimum height=5.8cm] (frame) at (5.4,3) {};
  \node[anchor=north,yshift=-2pt] at (frame.north) {\small\textbf{用户展示模块（全部功能点）}};
  \node (actor) at (-0.2,3) {\shortstack{\small$\odot$\\[-1pt]\scriptsize 普通用户\\[-2pt]\tiny （已登录）}};
  \node[uc] (u1) at (3.8,4.85) {个人信息展示\\FR-YHZS-0001};
  \node[uc] (u2) at (7,4.85) {检测记录展示\\FR-YHZS-0002};
  \node[uc] (u3) at (3.8,3.55) {审核任务展示\\FR-YHZS-0003};
  \node[uc] (u4) at (7,3.55) {社区反馈展示\\FR-YHZS-0004};
  \node[uc] (u5) at (3.8,2.25) {综合鉴伪报告\\FR-YHZS-0005};
  \node[uc] (u6) at (7,2.25) {人工审核进度\\与结果 FR-YHZS-0006};
  \foreach \x in {u1,u2,u3,u4,u5,u6} {
    \draw[arr] (actor.east) -- ++(0.22,0) |- (\x.west);
  }
\end{tikzpicture}
}
\end{center}

\subsubsection{功能点——个人信息展示（FR-YHZS-0001）}

\paragraph*{简述}

用户可通过该功能查看个人基本信息，包括昵称、ID、头像、等级等信息。

\paragraph*{前提条件}

用户已注册并登录系统，且系统数据库存储了该用户需要展示的信息。

\paragraph*{主要流程}

\begin{enumerate}
    \item 用户进入“个人中心”页面。
    \item 系统从数据库调取用户信息并渲染至前端页面。
    \item 用户可浏览个人基础信息，但不可直接修改（需通过其他模块操作）。
\end{enumerate}

\paragraph*{后继结果}

用户成功查看个人资料，页面显示完整信息。若数据加载失败，在该数据展示的位置提示“数据加载失败”。

\subsubsection{功能点——检测记录展示（FR-YHZS-0002）}

\paragraph*{简述}

用户可按时间顺序查看本人发起的学术检测任务记录，覆盖学术图像、全篇论文与同行评审 Review 等类型；支持筛选、检索、进入详情及下载对应检测报告。

\paragraph*{前提条件}

用户已登录；系统中已持久化该用户的检测任务数据。

\paragraph*{主要流程}

\begin{enumerate}
    \item 用户进入“检测记录”页面。
    \item 系统默认按时间倒序展示一定时间窗口内（如半年内，具体以产品配置为准）的检测任务条目。
    \item 用户可通过时间范围、内容类型等条件筛选，或通过关键字检索任务。
    \item 用户点击某条任务进入详情页，查看检测状态、子模块结果摘要及报告入口。
    \item 用户在详情或列表提供的入口下载检测报告（格式以后端约定为准，如 PDF）。
    \item 用户可对单条记录发起删除，经二次确认后由系统更新列表与存储。
\end{enumerate}

\paragraph*{后继结果}

列表与详情正确展示任务状态与可下载资源；无数据时展示空状态提示。删除成功后对应记录不再出现在列表中。加载或操作失败时，在界面给出明确错误提示并可重试。

\subsubsection{功能点——审核任务展示（FR-YHZS-0003）}

\paragraph*{简述}

用户可查看半年内提交的人工审核任务记录，支持查看学术图像、全篇论文及同行评审Review的人工审核详情、审核进度、下载阶段性审核报告及终止审核任务。

\paragraph*{前提条件}

用户已提交过人工审核请求，且对应审核任务数据已成功存储至系统数据库，包括图像检测任务、论文检测任务或Review检测任务。

\paragraph*{主要流程}

\begin{enumerate}
    \item 用户进入“人工审核”页面。
    \item 系统按时间倒序展示用户半年内提交的人工审核任务条目。
    \item 每个任务条目展示任务类型（图像 / 论文 / Review）、提交时间、当前审核状态、当前参与审核人数及审核进度。
    \item 用户点击任务条目后进入审核详情页面。
    \item 系统展示当前人工审核进度，包括已完成审核人数、当前审核意见汇总、AI初检结果及阶段性人工审核结论。
    \item 用户可下载当前阶段的综合审核报告，报告包括AI检测结果、人工审核意见汇总、风险等级及使用建议。
    \item 对于未完成的审核任务，用户可选择“终止审核”操作。
    \item 系统弹出二次确认窗口，确认后更新任务状态。
\end{enumerate}

\paragraph*{后继结果}

若存在人工审核任务，则用户成功查看审核任务状态、审核进度及阶段性结果；若无记录，则显示“还没有检测记录”。

终止操作完成后，任务状态更新为“已取消”，其他用户无法继续参与审核，系统同步生成终止记录并写入日志模块。

\subsubsection{功能点——社区反馈展示（FR-YHZS-0004）}

\paragraph*{简述}

用户可查看管理员对举报、评论及人工审核相关处理结果，以及个人审核任务、评论回复、论文检测结果反馈、Review检测反馈和系统消息通知。

\paragraph*{前提条件}

用户曾参与过评论、举报、人工审核或提交过检测任务。

若用户参与过举报操作，要求管理员已完成相关处理并更新数据；

若用户提交过论文、图像或Review检测任务，要求系统已生成对应反馈信息。

\paragraph*{主要流程}

\begin{enumerate}
    \item 用户进入“社区反馈”页面。
    \item 系统按时间顺序展示反馈消息列表。
    \item 反馈内容包括：
    \begin{enumerate}
        \item 管理员对举报内容的处理结果；
        \item 用户评论或审核意见的新回复；
        \item 人工审核结果更新通知；
        \item 论文AIGC检测结果更新通知；
        \item Review虚假评审检测反馈；
        \item 系统生成的综合鉴伪报告通知。
    \end{enumerate}
    \item 用户点击消息条目后查看详细反馈内容。
    \item 对于检测反馈消息，用户可直接跳转至对应检测记录详情页面。
\end{enumerate}

\paragraph*{后继结果}

若存在新的个人消息回复、举报处理结果更新、人工审核进度更新或检测结果更新，则标注为“新消息”；历史已读信息标注为“已读”。

若无任何反馈消息，则页面展示“暂无消息”。

\subsubsection{功能点——综合鉴伪报告下载与查看（FR-YHZS-0005）}

\paragraph*{简述}

用户针对某次检测任务，在同一入口在线查看系统生成的综合鉴伪报告，并支持将报告下载到本地。报告内容聚合该任务下各子检测维度（如论文 AIGC、图像可疑区域、Review 相关结论等，以实际任务类型与后端返回为准），便于用户一次性把握整体鉴伪结论与依据摘要。

\paragraph*{前提条件}

用户已登录；对应检测任务已存在且用户具备查看权限；综合报告已生成或系统支持展示“生成中/失败”等中间状态（以后端与产品约定为准）。

\paragraph*{主要流程}

\begin{enumerate}
    \item 用户在“检测记录”或任务详情中进入“综合鉴伪报告”入口。
    \item 前端请求报告元数据与正文结构（章节、图表、结论文本等），并渲染为可浏览页面。
    \item 用户可展开/折叠章节、定位到关键结论与引用数据说明（交互细节以实现为准）。
    \item 用户点击“下载”，获取与线上一致的报告文件（如 PDF）；若报告尚未就绪，展示预计状态或排队信息。
\end{enumerate}

\paragraph*{后继结果}

报告成功展示并可下载；未生成或无权限时给出对应提示。下载失败时提示原因并支持重试。

\paragraph*{补充说明}

报告字段与排版模板由前后端约定；若某类子检测未参与本次任务，报告中对应章节可隐藏或标注“不适用”。

\subsubsection{功能点——人工审核进度与结果查询（FR-YHZS-0006）}

\paragraph*{简述}

用户可查询本人发起的人工审核任务在处理流程中的当前进度、阶段性状态与最终结果，并支持按内容类型（如论文、Review 等）在统一列表或筛选视图中查看，与“检测记录”等模块形成互补入口。

\paragraph*{前提条件}

用户已登录；系统中存在与该用户关联的人工审核任务数据（含进行中、已完成、已取消等状态）。

\paragraph*{主要流程}

\begin{enumerate}
    \item 用户进入“人工审核进度/我的审核”等相关页面（名称以实现为准）。
    \item 系统按时间倒序或按状态分组展示任务列表，展示字段包括任务类型、提交时间、当前状态、参与人数或轮次等（以后端返回为准）。
    \item 用户可通过类型、状态、时间范围筛选或关键字搜索目标任务。
    \item 用户点击任务进入详情，查看进度时间线、各审核轮次摘要及平台汇总后的结论（若已产出）。
    \item 若业务支持，用户可在详情中查看或下载与人工审核相关的结果说明或附件。
\end{enumerate}

\paragraph*{后继结果}

用户能区分“待分配/审核中/已完成/已取消”等状态并查看终态说明；无任务时展示空状态。数据拉取失败时提示错误并可刷新重试。
